% plan.tex -- Example initial project plan / proposal using the qproposal class
%
% This template is intended for short project plans and proposals, such as
% initial plans for dissertation / project modules. It follows the structure
% suggested for typical initial plans (project description, aims and objectives,
% feasibility, work plan, references), but remains generic so it can be reused
% in other contexts.
%
% Recommended length for an initial plan of this style is around 1,500 words
% (excluding references, figures and tables), but this is *not* enforced by
% the class - adjust as needed for your institution or project.
%
% SPDX-FileCopyrightText: 2026 Frank Langbein <frank@langbein.org>
% SPDX-License-Identifier: LPPL-1.3c

\documentclass[sans,oneside,single,parskip,raggedright,numeric]{qproposal}

% ------------------------------------------------------------------------------
% Metadata
% ------------------------------------------------------------------------------

% The title will typically start with "Initial Plan:" followed by your project title.
\title{Initial Plan: A Sample Project Title}

% Your name (and optionally co-authors, if appropriate).
\author{Student Name}

% Optional supervisor name - printed under the author line when set.
\supervisor{Supervisor Name}

% Optional date (defaults to \today if omitted).
% \date{2 March 2026}

% Optional SPDX license identifier. If set, \bib will add a short License section
% before the references.
\license{CC-BY-NC-SA-4.0}
\copyrightholder{Student Name}
\licenseyear{2026}

% If you want keywords to appear at the end of the abstract, set them here.
\keywords{project planning, initial plan, dissertation}

% ------------------------------------------------------------------------------
% Document body
% ------------------------------------------------------------------------------

\begin{document}

\maketitle

% A short overview of the project and the purpose of the plan (typically
% a small paragraph). This does not count towards the usual word limits
% in many guidelines, but check your local rules.
\begin{abstract}
  This document outlines the initial plan for a sample project. It briefly
  describes the problem and its context, states aims and objectives, discusses
  feasibility issues, and presents a work plan with milestones and deliverables.
\end{abstract}

% ---------------------------------------------------------------------------
% 1. Project Description
% ---------------------------------------------------------------------------
%
% Briefly describe the project, its context and motivation, and the specific
% problem you intend to address. Explain the background only to the extent
% needed to understand the plan; a full literature review is not required here.
%
% You may also justify why the problem is an appropriate challenge for a
% final-year / master-level project, and clarify the scope (what is in and
% out of scope for this plan).

\section{Project Description}

Describe the overall context and motivation for your project, followed by a
clear and specific problem statement. Outline which aspects of the problem
you will address within the constraints of the project and which aspects
are explicitly out of scope~\citep{example-paper}. According
to~\citet{example-book}, we will address the problem in the following way.

% ---------------------------------------------------------------------------
% 2. Aims and Objectives
% ---------------------------------------------------------------------------
%
% State an overall aim (one or two sentences) and then list two to four
% specific objectives you intend to achieve. Each objective should be
% concrete and realistic for the available time, while still providing
% sufficient challenge.
%
% For each objective, briefly note any likely or highly likely risks and how
% you would mitigate them (for example, by scaling back scope or relaxing
% non-essential requirements). There is no need to consider unlikely risks,
% and personal extenuating circumstances are usually handled separately.

\section{Aims and Objectives}

The overall aim of this project is to \ldots

\begin{itemize}
  \item Objective 1: \ldots\ (include a brief note on likely risks and possible
    mitigations if relevant).
  \item Objective 2: \ldots
  \item Objective 3: \ldots
\end{itemize}

% ---------------------------------------------------------------------------
% 3. Feasibility
% ---------------------------------------------------------------------------
%
% Discuss any issues that may affect the feasibility of the project, such as:
%   - ethical approval (e.g. user studies, personal data),
%   - legal issues (e.g. intellectual property, licensing of software or data),
%   - special resources (e.g. non-standard hardware, paid APIs or datasets).
%
% For each issue, describe how it will be handled and what the impact on the
% project would be if things do not go as planned (risk and mitigation).
% If you reasonably expect no such issues, state this briefly with justification.

\section{Feasibility}

Summarise any ethical, legal or resource-related constraints relevant to
the project, and explain how you will address them. If you depend on external
approvals or resources, outline what happens if these are delayed or not
available and how you would adapt the project in that case.

% ---------------------------------------------------------------------------
% 4. Work Plan
% ---------------------------------------------------------------------------
%
% Provide a work plan over the relevant period (for example, weeks 1–12 of a
% term), including:
%   - clear milestones (what you expect to have achieved by a given date),
%   - tasks leading up to each milestone,
%   - expected deliverables (e.g. prototype, evaluation results, report draft).
%
% A simple week-by-week list with dates, milestones and main tasks is often
% sufficient. Consider other commitments (modules, exams, holidays) and reserve
% time for evaluation and writing towards the end of the project.

\section{Work Plan}

Outline your planned work over the project period. For example:

\begin{description}
  \item[Weeks 1--2] Refine project description, confirm feasibility, and agree
    aims and objectives with the supervisor. Identify key background reading.
  \item[Weeks 3--6] Design and implement the main solution approach in one or
    more iterations. Record design decisions and intermediate results.
  \item[Weeks 7--9] Evaluate the approach (e.g. experiments, user studies),
    analyse results, and, if time permits, refine or extend the solution.
  \item[Weeks 10--12] Consolidate documentation and write the main report
    chapters. Reserve buffer time for tasks that overrun and final polishing.
\end{description}

You may replace this list by a more detailed table or Gantt chart if you
prefer, as long as milestones, tasks and deliverables are clear.

% ---------------------------------------------------------------------------
% References
% ---------------------------------------------------------------------------
%
% If you cite any sources in the plan (recommended where appropriate), include
% them here. The qproposal class uses natbib; \bib will choose a numeric or
% author-year style based on the class options.

% Use \bib with an external .bib file (an example bibliography.bib is provided
% with the qdissertation package).
\bib{bibliography}

\end{document}
